% Vietnamese translation for OI Public Library
% Source-PDF: 比赛 CONTEST/2021“MINIEYE杯”中国大学生算法设计超级联赛/2021“MINIEYE杯”中国大学生算法设计超级联赛（9）/2021“MINIEYE杯”中国大学生算法设计超级联赛（9）-题目集.pdf
\documentclass[11pt]{article}
\usepackage{oipl-vn}

\title{Đề bài HDU Multi-University 2021, Contest 9}
\author{}
\date{}

\begin{document}
\maketitle

\origtitle{2021 HDU Multi-University Training Contest 9}
\sourcepath{contest/hdu-2021-minieye-09-problems.pdf}

\section{Problem A. NJU Emulator}

\paragraph{Tệp vào:} standard input
\paragraph{Tệp ra:} standard output
\paragraph{Giới hạn thời gian:} 1 second
\paragraph{Giới hạn bộ nhớ:} 256 megabytes

``Dr.JYY, tập lệnh s86, và NEMU, chúng làm bạn nhớ đến điều gì?''

``Ôi, không...''

NJU Emulator, còn gọi là NEMU, là trình mô phỏng kiến trúc s86 mới nhất do
Dr. JYY phát triển. s86 là một kiến trúc máy tính kiểu ngăn xếp, trong đó các
lệnh máy chỉ thao tác trên phần tử ở đỉnh ngăn xếp.

Mô hình tính toán của s86 gồm một ngăn xếp và một chương trình có độ dài hữu
hạn. Mỗi phần tử trong ngăn xếp là một số nguyên không dấu 64 bit. Ký hiệu
ngăn xếp là \(S\), kích thước ngăn xếp là \(\operatorname{size}(S)\). Chương
trình gồm các lệnh trong bảng sau và phải kết thúc bằng một lệnh dừng.

Khi máy s86 chạy, trước hết ngăn xếp được khởi tạo rỗng, sau đó các lệnh trong
chương trình được thực hiện lần lượt cho đến lệnh cuối cùng. Sau một lệnh
\texttt{end}, máy sẽ in phần tử ở đỉnh ngăn xếp và dừng.

\begin{center}
\begin{tabular}{lll}
\textbf{Ký hiệu} & \textbf{Chức năng} & \textbf{Ràng buộc}\\
\texttt{p1} & Đưa hằng số \(1\) vào \(S\) & Không có ràng buộc\\
\texttt{dup} & Thêm một bản sao của phần tử đỉnh vào \(S\) &
  \(\operatorname{size}(S)\ge 1\)\\
\texttt{pop} & Lấy phần tử đỉnh ra khỏi \(S\) &
  \(\operatorname{size}(S)\ge 1\)\\
\texttt{swap} & Đổi chỗ hai phần tử trên cùng của \(S\) &
  \(\operatorname{size}(S)\ge 2\)\\
\texttt{add x} & Cộng phần tử thứ \(x\) vào phần tử đỉnh của \(S\) &
  \(x\ge 0\) và \(\operatorname{size}(S)>x\)\\
\texttt{sub x} & Trừ phần tử thứ \(x\) khỏi phần tử đỉnh của \(S\) &
  \(x\ge 0\) và \(\operatorname{size}(S)>x\)\\
\texttt{mul x} & Nhân phần tử thứ \(x\) với phần tử đỉnh của \(S\) &
  \(x\ge 0\) và \(\operatorname{size}(S)>x\)\\
\texttt{end} & In phần tử đỉnh của \(S\) và dừng chương trình &
  \(\operatorname{size}(S)\ge 1\)
\end{tabular}
\end{center}

Trong bảng, phần tử thứ \(x\) là phần tử thứ \(x\) tính từ trên xuống dưới
của ngăn xếp, trong đó phần tử đỉnh là phần tử thứ \(0\).

Đáng chú ý, mọi phép toán số học \texttt{add}, \texttt{sub}, \texttt{mul}
trong tập lệnh s86 phải được thực hiện theo modulo \(2^{64}\). Nghĩa là nếu
kết quả toán học của phép toán là \(X\), thì kết quả trong tập lệnh s86 là
\(X'\) với \(0\le X'<2^{64}\) và \(X-X'\) là bội của \(2^{64}\). Có thể chứng
minh rằng \(X'\) như vậy tồn tại với mọi số nguyên \(X\).

Sau khi hoàn thành NEMU, để kiểm tra tính đúng đắn của nó, Dr. JYY muốn bạn
viết một chương trình bằng các lệnh s86 sao cho khi chương trình dừng, nó in
ra số nguyên \(N\) đã cho.

\subsection*{Dữ liệu vào}

Dòng đầu tiên chứa một số \(T\) (\(1\le T\le 10000\)), là số bộ dữ liệu.

Mỗi bộ dữ liệu chứa một số nguyên \(N\) (\(0\le N<2^{64}\)) trên một dòng, là
số mà chương trình của bạn cần in ra khi dừng.

\subsection*{Dữ liệu ra}

Với mỗi bộ dữ liệu, hãy in một đoạn chương trình gồm các lệnh s86 ở đúng dạng
như trong bảng của đề bài, sao cho khi chương trình dừng, nó in ra số nguyên
\(N\) đã cho.

Ngoài ra, chương trình của bạn phải chứa không quá \(50\) lệnh s86. Bảo đảm
rằng tồn tại ít nhất một chương trình như vậy trong các giới hạn của bài. Nếu
bất kỳ ràng buộc nào trong bảng lệnh bị vi phạm trong khi thực thi chương
trình, lời giải của bạn sẽ bị xem là sai.

Nếu có nhiều chương trình thỏa mãn ràng buộc, bạn có thể in ra bất kỳ chương
trình nào trong số đó.

\subsection*{Ví dụ}

\paragraph{standard input}
\begin{verbatim}
3
0
2
18446744073709551615
\end{verbatim}

\paragraph{standard output}
\begin{verbatim}
p1
sub 0
end
p1
add 0
end
p1
dup
add 1
swap
sub 1
end
\end{verbatim}

\subsection*{Ghi chú}

Quá trình thực thi bộ dữ liệu cuối cùng trong output mẫu được trình bày trong
bảng sau.

\begin{center}
\begin{tabular}{ll}
\textbf{Lệnh} & \textbf{Ngăn xếp}\\
\texttt{p1} & \([1)\)\\
\texttt{dup} & \([1,1)\)\\
\texttt{add 1} & \([1,2)\)\\
\texttt{swap} & \([2,1)\)\\
\texttt{sub 1} & \([2,18446744073709551615)\)\\
\texttt{end} & Chương trình in \(18446744073709551615\) và dừng
\end{tabular}
\end{center}

\section{Problem B. Just another board game}

\paragraph{Tệp vào:} standard input
\paragraph{Tệp ra:} standard output
\paragraph{Giới hạn thời gian:} 1 second
\paragraph{Giới hạn bộ nhớ:} 256 megabytes

``Vậy giờ tôi đi quân đến \((179,231)\). Đây là lượt thứ \(999999999\) của
ván này. Cuối cùng, chỉ còn một lượt nữa!''

``Gì cơ? Chẳng phải đây mới là lượt thứ \(999999997\) sao?''

``Ôi, chết tiệt.''

Sau khi chơi vài ván cờ vây, Roundgod và kimoyami quyết định thử một trò chơi
khác. Giờ họ chơi một loại trò chơi mới trên một bàn cờ dạng lưới có \(n\)
hàng và \(m\) cột. Giả sử ô góc trên bên trái có tọa độ \((1,1)\), còn ô góc
dưới bên phải có tọa độ \((n,m)\). Mỗi ô trên bàn có một số; số viết trên ô ở
hàng \(i\), cột \(j\) là \(a_{ij}\). Ban đầu có một quân cờ ở góc trên bên
trái, tức là ô \((1,1)\). Hai người chơi luân phiên chọn một trong các thao
tác sau, bắt đầu từ Roundgod:
\begin{enumerate}
  \item Di chuyển quân cờ. Nếu đến lượt Roundgod, anh ta có thể di chuyển quân
  cờ đến bất kỳ vị trí nào trên cùng hàng. Cũng được phép di chuyển đến vị trí
  hiện tại, tức là không di chuyển. Nếu đến lượt kimoyami, anh ta có thể di
  chuyển quân cờ đến bất kỳ vị trí nào trên cùng cột. Cũng được phép không di
  chuyển.
  \item Bỏ cuộc và về nhà. Kết thúc trò chơi ngay lập tức.
\end{enumerate}

Trò chơi kết thúc khi một trong hai người chọn thao tác thứ hai hoặc khi trò
chơi đã diễn ra \(k\) lượt. Mỗi thao tác của một trong hai người đều được tính
là một lượt. Giá trị của trò chơi được định nghĩa là số trên ô mà quân cờ đang
đứng khi trò chơi kết thúc.

Roundgod muốn tối đa hóa giá trị này, còn kimoyami muốn tối thiểu hóa nó. Họ
không đủ kiên nhẫn để thật sự chơi một ván có thể kéo dài như vậy, nên muốn
bạn tính giá trị cuối cùng của trò chơi nếu cả hai đều chọn chiến lược tối ưu.

\subsection*{Dữ liệu vào}

Dòng đầu tiên chứa một số \(T\) (\(1\le T\le 25\)), là số bộ dữ liệu.

Dòng đầu tiên của mỗi bộ dữ liệu chứa ba số nguyên \(n,m,k\)
\[
  n,m\ge 1,\qquad 1\le n\times m\le 10^5,\qquad 1\le k\le 10^{18},
\]
lần lượt là kích thước bàn cờ và số lượt tối đa mà trò chơi sẽ kéo dài.

Sau đó là \(n\) dòng. Dòng thứ \(i\) (\(1\le i\le n\)) chứa \(m\) số nguyên
\(a_{i1},a_{i2},\ldots,a_{im}\), trong đó \(a_{ij}\)
\((0\le a_{ij}\le 10^9)\) là số viết trên ô ở hàng \(i\), cột \(j\).

Bảo đảm rằng \(\sum(n\times m)\le 10^6\) trên tất cả các bộ dữ liệu.

\subsection*{Dữ liệu ra}

Với mỗi bộ dữ liệu, in một số nguyên trên một dòng, là giá trị cuối cùng của
trò chơi nếu cả hai người chơi đều chọn chiến lược tối ưu.

\subsection*{Ví dụ}

\paragraph{standard input}
\begin{verbatim}
3
2 2 2
1 2
2 1
2 2 1
1 2
2 1
2 3 2
1 3 2
3 2 1
\end{verbatim}

\paragraph{standard output}
\begin{verbatim}
1
2
2
\end{verbatim}

\section{Problem C. Dota2 Pro Circuit}

\paragraph{Tệp vào:} standard input
\paragraph{Tệp ra:} standard output
\paragraph{Giới hạn thời gian:} 2 seconds
\paragraph{Giới hạn bộ nhớ:} 256 megabytes

TI10 và ICPC World Finals 2020, sự kiện nào sẽ được tổ chức trước? Hãy cá
cược nào!

The International (TI) là sự kiện lớn nhất và danh giá nhất của Dota2, thường
được tổ chức hằng năm vào tháng Tám. Các đội Dota2 cần cố gắng kiếm điểm để
đủ điều kiện tham dự TI. Có hai cách kiếm điểm: thứ nhất là thi đấu ở các giải
khu vực, thứ hai là thi đấu ở các giải đấu nơi tất cả các đội tụ họp và kiếm
điểm theo thứ hạng trong giải. Điểm cuối cùng của một đội là tổng điểm từ cả
giải khu vực và giải đấu.

Hiện các giải khu vực đã kết thúc. Có \(n\) đội tham gia, và đội thứ \(i\) đã
kiếm được \(a_i\) điểm từ các giải khu vực. Ngoài ra, đội đạt hạng \(i\) trong
giải đấu sẽ nhận được \(b_i\) điểm.

cyz là một người hâm mộ Dota2 cuồng nhiệt. Vì vậy trước khi giải đấu bắt đầu,
anh ấy sẽ dự đoán thứ hạng cuối cùng của tất cả các đội. cyz muốn biết, với
mỗi đội, thứ hạng tốt nhất có thể và thứ hạng tệ nhất có thể của đội đó sau
khi giải đấu kết thúc.

Nếu một đội có điểm cuối cùng bằng \(x\), thứ hạng của đội đó được định nghĩa
là \(1\) cộng với số đội có điểm lớn hơn hẳn \(x\). Ví dụ, nếu điểm cuối cùng
của bốn đội lần lượt là \(700,500,500,300\), thì thứ hạng cuối cùng của họ lần
lượt là \(1,2,2,4\).

\subsection*{Dữ liệu vào}

Dòng đầu tiên chứa một số \(T\) (\(1\le T\le 20\)), là số bộ dữ liệu.

Dòng đầu tiên của mỗi bộ dữ liệu chứa một số \(n\) (\(1\le n\le 5000\)), là
số đội khác nhau tham gia các giải khu vực và giải đấu.

Dòng tiếp theo chứa \(n\) số nguyên \(a_1,a_2,\ldots,a_n\)
\((0\le a_i\le 10^9)\), là điểm của từng đội trước khi giải đấu bắt đầu.

Sau đó là một dòng chứa \(n\) số nguyên \(b_1,b_2,\ldots,b_n\)
\[
  0\le b_n\le b_{n-1}\le\cdots\le b_1\le 10^9,
\]
trong đó \(b_i\) là số điểm mà một đội nhận được nếu xếp hạng \(i\) trong
giải đấu.

Bảo đảm rằng có không quá \(8\) bộ dữ liệu có \(n>100\).

\subsection*{Dữ liệu ra}

Với mỗi bộ dữ liệu, in \(n\) dòng. Dòng thứ \(i\) chứa hai số nguyên
\(\operatorname{best}_i,\operatorname{worst}_i\)
\[
  1\le \operatorname{best}_i\le \operatorname{worst}_i\le n,
\]
lần lượt là thứ hạng tốt nhất và tệ nhất có thể của đội thứ \(i\) sau khi
giải đấu kết thúc.

\subsection*{Ví dụ}

\paragraph{standard input}
\begin{verbatim}
2
3
5 10 8
5 2 1
2
5 6
4 4
\end{verbatim}

\paragraph{standard output}
\begin{verbatim}
2 3
1 2
1 3
2 2
1 1
\end{verbatim}

\section{Problem D. Into the Woods}

\paragraph{Tệp vào:} standard input
\paragraph{Tệp ra:} standard output
\paragraph{Giới hạn thời gian:} 3 seconds
\paragraph{Giới hạn bộ nhớ:} 256 megabytes

Không phải mọi người lang thang đều bị lạc.

Có lẽ chỉ là một cuộc đi dạo, hoặc có lẽ là để trốn khỏi điều gì đó. Roundgod
tạm thời rời khỏi thành phố nơi anh từng sống và tiến vào một khu rừng vô
tận.

Roundgod bắt đầu lang thang trong rừng. Nếu xem khu rừng là một mặt phẳng
tọa độ Descartes hai chiều, thì Roundgod bắt đầu tại \((0,0)\). Mỗi lần, anh
chọn ngẫu nhiên đều một trong bốn hướng và đi thêm một bước, tức là các thay
đổi tọa độ có thể là \((-1,0)\), \((1,0)\), \((0,-1)\), \((0,1)\).

Sau khi đi \(n\) bước, anh dừng lại trong trạng thái mơ hồ, nhận ra rằng mình
đã lang thang quá lâu và đã đến lúc quay về cuộc sống thường ngày. Giờ chỉ
còn một câu hỏi khiến anh băn khoăn: trong \(n\) bước đã đi, anh đã từng cách
điểm xuất phát bao xa? Vì Roundgod đã quên các bước mình đã đi, anh chỉ muốn
biết kỳ vọng của câu trả lời trên tất cả các lộ trình có thể.

Nói hình thức hơn, Roundgod bắt đầu tại \((0,0)\), mỗi lần chọn ngẫu nhiên đều
một trong bốn hướng và đi thêm một bước. Bạn cần tính, trong tất cả \(4^n\) lộ
trình có thể, kỳ vọng của khoảng cách Manhattan lớn nhất giữa bất kỳ vị trí
nào trên lộ trình và \((0,0)\).

Cụ thể, khoảng cách Manhattan giữa hai điểm \(p_1=(x_1,y_1)\) và
\(p_2=(x_2,y_2)\) trên mặt phẳng Descartes được định nghĩa là
\[
  d(p_1,p_2)=|x_1-x_2|+|y_1-y_2|.
\]
Ngoài ra, với mọi modulo nguyên tố \(10^8\le MOD\le 10^9+7\), có thể chứng
minh rằng đáp án viết được dưới dạng phân số \(\frac{P}{Q}\), trong đó
\(P,Q\) nguyên tố cùng nhau và \(Q\not\equiv 0\pmod {MOD}\). Bạn cần in
\[
  P\cdot Q^{-1}\pmod {MOD},
\]
trong đó \(Q^{-1}\pmod {MOD}\) là nghịch đảo modulo của \(Q\) theo \(MOD\).

\subsection*{Dữ liệu vào}

Dòng đầu tiên chứa một số \(T\) (\(1\le T\le 10\)), là số bộ dữ liệu.

Dòng đầu tiên của mỗi bộ dữ liệu chứa hai số \(n,MOD\)
\[
  1\le n\le 10^6,\qquad 10^8\le MOD\le 10^9+7,
\]
lần lượt là số bước đã đi và modulo. Bảo đảm rằng \(MOD\) là số nguyên tố.

\subsection*{Dữ liệu ra}

Với mỗi bộ dữ liệu, in một số nguyên trên một dòng, là đáp án.

\subsection*{Ví dụ}

\paragraph{standard input}
\begin{verbatim}
3
1 1000000007
2 998244353
50 1000000007
\end{verbatim}

\paragraph{standard output}
\begin{verbatim}
1
249561090
603242629
\end{verbatim}

\section{Problem E. Did I miss the lethal?}

\paragraph{Tệp vào:} standard input
\paragraph{Tệp ra:} standard output
\paragraph{Giới hạn thời gian:} 6 seconds
\paragraph{Giới hạn bộ nhớ:} 256 megabytes

Bạn có thể gây bao nhiêu sát thương với hai lá ``Soulfire'' trên tay? \(8\)?
Bạn chắc chứ?

-skyline- đang chơi Hearthstone, và lớp nhân vật yêu thích của anh là
Warlock. Có một cơ chế đặc biệt được thiết kế cho các lá bài của lớp Warlock:
sau khi đánh một số lá từ tay, người chơi có thể phải bỏ ngẫu nhiên một số lá
khác còn lại trên tay.

-skyline- có \(n\) lá bài trên tay. Mỗi lá có thể gây một lượng sát thương
nào đó, nhưng cũng có thể khiến anh phải bỏ một số lá khác trên tay sau khi
đánh. Anh đang tính toán một lượt kết liễu. Là một người chơi Hearthstone
chuyên nghiệp, anh không tin vào vận may của mình, nên muốn tính trong trường
hợp xấu nhất, lượng sát thương tối đa mà anh có thể gây ra bằng tay bài này là
bao nhiêu. Bỏ qua chi phí mana của các lá bài và giả sử mọi lá đều có thể
được đánh nếu chưa bị bỏ.

Nói hình thức hơn, có \(n\) lá bài trên tay của -skyline-. Lá thứ \(i\) có hai
thuộc tính \(d_i\) và \(a_i\), lần lượt là sát thương nó gây ra và số lá cần
bỏ sau khi đánh lá này. -skyline- có thể chọn thứ tự đánh các lá từ tay, từng
lá một. Sau khi anh chọn đánh lá thứ \(i\) từ tay, lá này không còn trên tay
nữa, anh gây \(d_i\) sát thương lên tướng địch, rồi \(a_i\) lá trong phần tay
còn lại được chọn ngẫu nhiên đều và bị bỏ. Nếu lúc đó -skyline- còn ít hơn
\(a_i\) lá trên tay, thì tất cả các lá còn lại đều bị bỏ. Sau khi một lá bị bỏ
khỏi tay, nó không thể được đánh nữa.

-skyline- muốn chọn chiến lược tốt nhất sao cho trong trường hợp xấu nhất, anh
gây được tổng sát thương lớn nhất lên tướng địch. Hãy giúp -skyline- tìm tổng
sát thương đó.

\subsection*{Dữ liệu vào}

Dòng đầu tiên chứa một số \(T\) (\(1\le T\le 20\)), là số bộ dữ liệu.

Dòng đầu tiên của mỗi bộ dữ liệu chứa một số \(n\) (\(1\le n\le 200\)), là số
lá bài trên tay của -skyline-.

Sau đó là \(n\) dòng. Dòng thứ \(i\) chứa hai số nguyên \(d_i\)
\((1\le d_i\le 10^7)\) và \(a_i\) (\(1\le a_i\le 4\)), lần lượt là sát thương
lá thứ \(i\) có thể gây ra và số lá cần bỏ sau khi đánh lá đó.

Bảo đảm rằng \(\sum n\le 2000\) trên tất cả các bộ dữ liệu.

\subsection*{Dữ liệu ra}

Với mỗi bộ dữ liệu, in một số nguyên trên một dòng, là tổng sát thương lớn
nhất mà -skyline- có thể gây lên tướng địch trong trường hợp xấu nhất.

\subsection*{Ví dụ}

\paragraph{standard input}
\begin{verbatim}
2
3
3 1
4 2
5 2
6
10 1
10 1
15 2
20 3
20 3
25 4
\end{verbatim}

\paragraph{standard output}
\begin{verbatim}
7
40
\end{verbatim}

\subsection*{Ghi chú}

Với bộ dữ liệu mẫu đầu tiên, chiến lược tối ưu của -skyline- là đánh lá thứ
nhất, sau đó đánh lá duy nhất còn lại chưa bị bỏ. Trong trường hợp này,
-skyline- có thể gây tổng cộng \(3+4=7\) sát thương lên tướng địch trong
trường hợp xấu nhất.

\section{Problem F. Guess The Weight}

\paragraph{Tệp vào:} standard input
\paragraph{Tệp ra:} standard output
\paragraph{Giới hạn thời gian:} 2 seconds
\paragraph{Giới hạn bộ nhớ:} 256 megabytes

``Guess The Weight, lượt 2. Được, hãy đánh lá này xem ta rút được gì.''

Rút được ``Innervate''.

``May quá! Mình chỉ cần bấm `Greater' là được rút thêm miễn phí!''

Thấy một lá ``Innervate'' nữa.

``Bạn có chắc muốn gỡ Hearthstone khỏi máy tính không? Có.''

uuzlovetree thích chơi Hearthstone, và lớp nhân vật yêu thích của anh là
Druid. Trong Hearthstone, Druid có một lá phép tên là ``Guess the weight''.

uuzlovetree biết số lá trong bộ bài và biết chi phí mana của từng lá. Anh
muốn biết xác suất rút được lá thứ hai khi đánh ``Guess the weight'' dưới
chiến lược đoán tối ưu.

Nói hình thức hơn, giả sử hiện có \(m\) lá trong bộ bài, mỗi lá có một số biểu
thị chi phí mana. Ta xáo ngẫu nhiên tất cả \(m\) lá trong bộ bài, nghĩa là mọi
cách sắp xếp trong \(m!\) cách đều có xác suất như nhau. Khi một người đánh
lá ``Guess the weight'', người đó rút lá đầu tiên của bộ bài và chọn một trong
hai lựa chọn sau:
\begin{enumerate}
  \item Dự đoán rằng lá tiếp theo, tức lá thứ hai của bộ bài, có chi phí mana
  nhỏ hơn lá đầu tiên; sau đó lật lá tiếp theo. Nếu dự đoán đúng, người đó rút
  lá thứ hai.
  \item Dự đoán rằng lá tiếp theo, tức lá thứ hai của bộ bài, có chi phí mana
  lớn hơn lá đầu tiên; sau đó lật lá tiếp theo. Nếu dự đoán đúng, người đó rút
  lá thứ hai.
\end{enumerate}

Chú ý: nếu lá thứ hai của bộ bài có chi phí mana bằng lá đầu tiên, thì dù bạn
chọn lựa chọn nào, bạn cũng không thể rút lá thứ hai.

Ban đầu có \(n\ge 2\) lá trong bộ bài, với chi phí mana lần lượt là
\(a_1,a_2,\ldots,a_n\). Sau đó có \(q\) sự kiện xảy ra với bộ bài. Mỗi sự
kiện thuộc một trong hai dạng sau:
\begin{enumerate}
  \item Thêm \(x\) lá có chi phí mana \(w\) vào bộ bài.
  \item Hỏi, khi áp dụng chiến lược tối ưu, xác suất lớn nhất để rút được lá
  thứ hai khi đánh ``Guess the weight'' là bao nhiêu.
\end{enumerate}

\subsection*{Dữ liệu vào}

Dòng đầu tiên chứa một số \(T\) (\(1\le T\le 10\)), là số bộ dữ liệu.

Dòng đầu tiên của mỗi bộ dữ liệu chứa hai số \(n\) (\(2\le n\le 2\cdot 10^5\))
và \(q\) (\(1\le q\le 2\cdot 10^5\)), lần lượt là kích thước ban đầu của bộ
bài và số sự kiện.

Sau đó là một dòng chứa \(n\) số nguyên \(a_1,a_2,\ldots,a_n\)
\((1\le a_i\le 2\cdot 10^5)\), là chi phí mana của \(n\) lá ban đầu.

Tiếp theo là \(q\) dòng, mỗi dòng thuộc một trong hai dạng:
\begin{enumerate}
  \item \texttt{1 x w}, nghĩa là thêm \(x\) lá có chi phí mana \(w\) vào bộ
  bài \((1\le x\le 100,\ 1\le w\le 2\cdot 10^5)\).
  \item \texttt{2}, nghĩa là một truy vấn hỏi khi đánh ``Guess the weight''
  với bộ bài hiện tại, xác suất lớn nhất để rút được lá thứ hai là bao nhiêu.
\end{enumerate}

Bảo đảm rằng \(\sum n\le 10^6\) và \(\sum q\le 10^6\) trên tất cả các bộ dữ
liệu.

\subsection*{Dữ liệu ra}

Với mỗi sự kiện loại hai của mỗi bộ dữ liệu, in một phân số dạng
\(\texttt{p/q}\) trên một dòng, là xác suất lớn nhất để rút được lá thứ hai.
Yêu cầu \(p\) và \(q\) đều là số nguyên và \(\gcd(p,q)=1\), trong đó
\(\gcd(p,q)\) là ước chung lớn nhất của \(p\) và \(q\).

\subsection*{Ví dụ}

\paragraph{standard input}
\begin{verbatim}
2
2 5
1 1
2
1 1 2
2
1 1 2
2
2 5
1 2
2
1 1 3
2
1 100 4
2
\end{verbatim}

\paragraph{standard output}
\begin{verbatim}
0/1
2/3
2/3
1/1
5/6
201/3502
\end{verbatim}

\subsection*{Ghi chú}

Với bộ dữ liệu mẫu đầu tiên, ban đầu bộ bài gồm hai lá có chi phí mana \(1\).
Vì vậy dù chọn cách nào, người chơi cũng không thể rút lá thứ hai.

Sau khi thêm một lá có chi phí mana \(2\) vào bộ bài, chiến lược tối ưu là như
sau: nếu chi phí mana của lá đầu tiên rút được là \(1\), dự đoán rằng lá tiếp
theo có chi phí lớn hơn; ngược lại, dự đoán rằng lá tiếp theo có chi phí nhỏ
hơn. Có thể kiểm chứng rằng với chiến lược này, xác suất rút được lá thứ hai
là \(\frac{2}{3}\).

Sau khi thêm một lá nữa có chi phí mana \(2\) vào bộ bài, chiến lược tối ưu
không thay đổi. Có thể kiểm chứng rằng với chiến lược này, xác suất rút được
lá thứ hai vẫn là \(\frac{2}{3}\).

\section{Problem G. Boring data structure problem}

\paragraph{Tệp vào:} standard input
\paragraph{Tệp ra:} standard output
\paragraph{Giới hạn thời gian:} 3 seconds
\paragraph{Giới hạn bộ nhớ:} 256 megabytes

``Sao bạn không viết một câu chuyện nền cho bài cấu trúc dữ liệu nhàm chán
này?''

``Vì nó quá nhàm chán...''

Có một hàng đợi, chính xác hơn là một hàng đợi hai đầu, ban đầu rỗng. Sau đó,
có vô hạn phần tử được đánh số \(1,2,\ldots\) đi vào hàng đợi theo thứ tự số
của chúng, tức là phần tử đầu tiên đi vào hàng đợi được đánh số \(1\), phần tử
thứ hai được đánh số \(2\), v.v. Nếu một phần tử rời khỏi hàng đợi, nó sẽ
không vào hàng đợi nữa.

Có \(q\) thao tác, mỗi thao tác thuộc một trong các dạng sau:
\begin{enumerate}
  \item Cho phần tử tiếp theo đi vào đầu trái của hàng đợi.
  \item Cho phần tử tiếp theo đi vào đầu phải của hàng đợi.
  \item Cho phần tử được đánh số \(x\) rời khỏi hàng đợi.
  \item Hỏi số của phần tử ở giữa hàng đợi. Nếu hiện có \(m\) phần tử trong
  hàng đợi, bạn cần in số của phần tử thứ \(\left\lceil\frac{m+1}{2}\right\rceil\)
  tính từ bên trái.
\end{enumerate}

\subsection*{Dữ liệu vào}

Dòng đầu tiên chứa một số \(q\) (\(1\le q\le 10^7\)), là số thao tác.

Sau đó là \(q\) dòng, mỗi dòng thuộc một trong các dạng sau:
\begin{enumerate}
  \item \texttt{L}, nghĩa là phần tử tiếp theo đi vào đầu trái của hàng đợi.
  \item \texttt{R}, nghĩa là phần tử tiếp theo đi vào đầu phải của hàng đợi.
  \item \texttt{G x}, nghĩa là phần tử được đánh số \(x\) rời khỏi hàng đợi.
  Bảo đảm rằng phần tử được đánh số \(x\) đang ở trong hàng đợi khi thao tác
  này được áp dụng.
  \item \texttt{Q}, nghĩa là một truy vấn hỏi số của phần tử ở giữa hàng đợi.
  Bảo đảm rằng hàng đợi không rỗng khi thao tác này được áp dụng.
\end{enumerate}

Bảo đảm rằng số thao tác loại ba và số thao tác loại bốn đều nhỏ hơn
\(1.5\cdot 10^6\).

\subsection*{Dữ liệu ra}

Với mỗi thao tác loại bốn, in một số trên một dòng, là đáp án cho truy vấn.

\subsection*{Ví dụ}

\paragraph{standard input}
\begin{verbatim}
9
L
L
L
Q
R
Q
G 1
R
Q
\end{verbatim}

\paragraph{standard output}
\begin{verbatim}
2
1
4
\end{verbatim}

\section{Problem H. Integers Have Friends 2.0}

\paragraph{Tệp vào:} standard input
\paragraph{Tệp ra:} standard output
\paragraph{Giới hạn thời gian:} 3 seconds
\paragraph{Giới hạn bộ nhớ:} 256 megabytes

Lời cảm ơn: đặc biệt cảm ơn bài Codeforces 1548E Integer Have Friends đã cung
cấp phát biểu tổng quát cho bài này.

Nhà toán học Ấn Độ Srinivasa Ramanujan từng trích dẫn câu nói nổi tiếng của
nhà toán học Ấn Độ Srinivasa Ramanujan(?) rằng ``mọi số nguyên dương đều là
một người bạn riêng của ông''.

Hóa ra các số nguyên dương cũng có thể là bạn của nhau! Bạn được cho một mảng
\(a\) gồm các số nguyên dương đôi một phân biệt.

Định nghĩa một dãy con \(a_{c_1},a_{c_2},\ldots,a_{c_k}\), trong đó \(k\ge 1\)
và \(1\le c_1<c_2<\cdots<c_k\le n\), là một nhóm bạn khi và chỉ khi tồn tại
một số nguyên \(m\ge 2\) sao cho
\[
  a_{c_1}\bmod m=a_{c_2}\bmod m=\cdots=a_{c_k}\bmod m,
\]
trong đó \(x\bmod y\) là số dư khi chia \(x\) cho \(y\).

Người bạn gispzjz của bạn muốn biết kích thước của nhóm bạn lớn nhất trong
\(a\). Bạn có thể giúp anh ấy không?

\subsection*{Dữ liệu vào}

Dòng đầu tiên chứa một số \(T\) (\(1\le T\le 30\)), là số bộ dữ liệu.

Dòng đầu tiên của mỗi bộ dữ liệu chứa một số nguyên \(n\)
\((2\le n\le 2\times 10^5)\), là kích thước của mảng \(a\).

Sau đó là một dòng chứa \(n\) số nguyên \(a_1,a_2,\ldots,a_n\)
\((1\le a_i\le 4\times 10^{12})\), biểu thị nội dung của mảng \(a\). Bảo đảm
rằng tất cả các số trong \(a\) đều phân biệt.

Bảo đảm rằng \(\sum n\le 10^6\) trên tất cả các bộ dữ liệu.

\subsection*{Dữ liệu ra}

Với mỗi bộ dữ liệu, in một dòng chứa một số nguyên duy nhất, là kích thước của
nhóm bạn lớn nhất trong \(a\).

\subsection*{Ví dụ}

\paragraph{standard input}
\begin{verbatim}
3
3
10 12 15
4
4 6 9 19
6
2 8 11 15 19 38
\end{verbatim}

\paragraph{standard output}
\begin{verbatim}
2
3
4
\end{verbatim}

\section{Problem I. Little Prince and the garden of roses}

\paragraph{Tệp vào:} standard input
\paragraph{Tệp ra:} standard output
\paragraph{Giới hạn thời gian:} 1 second
\paragraph{Giới hạn bộ nhớ:} 256 megabytes

Chàng hoàng tử bé đứng trước một khu vườn, nơi hoa hồng đang nở rộ.

Giả sử vườn hồng có \(n\) hàng hoa hồng, mỗi hàng có \(n\) bông, xếp thành
một hình vuông. Ban đầu, mỗi bông hồng có một màu, dĩ nhiên là màu của cánh
hoa. Bông hồng ở hàng thứ \(i\), cột thứ \(j\) có màu \(c_{ij}\). Chắc chắn
có những bông hồng đủ màu sắc; sao mọi người cứ nghĩ đến màu đỏ khi nói về
hoa hồng?

Khi hoàng tử bé đi dạo trong vườn hồng, cậu sẽ bật khóc nếu nhìn thấy hai
bông hồng có cùng màu trong cùng một hàng hoặc cùng một cột, vì điều đó làm
cậu nhớ rằng bông hoa của mình không phải là duy nhất trên thế giới.

Bạn, một người tốt bụng, chắc chắn không muốn thấy hoàng tử bé khóc. May mắn
thay, bạn đến vườn hồng ngay trước cậu. Bạn có một cây cọ thần kỳ có thể sơn
thân của một số bông hồng thành các màu khác nhau; bạn đâu muốn làm hỏng cánh
hoa của chúng. Ban đầu, thân của tất cả hoa hồng đều màu xanh lá. Bạn muốn
sơn thân của một số bông hồng thành các màu khác nhau, sao cho không còn hai
bông hồng nào trong cùng một hàng hoặc cùng một cột có cả màu cánh hoa và màu
thân giống nhau.

Vấn đề duy nhất là số lượng loại màu khác nhau bạn cần chuẩn bị. Vì việc chuẩn
bị các loại màu khác nhau tốn thời gian, bạn muốn dùng ít loại màu nhất có
thể. Hãy tìm một cách sơn hoa sao cho điều kiện được thỏa mãn. Miễn là số
lượng loại màu được tối thiểu hóa, mọi cách sơn hợp lệ đều được chấp nhận.

\subsection*{Dữ liệu vào}

Dòng đầu tiên chứa một số \(T\) (\(1\le T\le 100\)), là số bộ dữ liệu.

Dòng đầu tiên của mỗi bộ dữ liệu chứa một số nguyên \(n\) (\(1\le n\le 300\)),
là kích thước của khu vườn. Sau đó là \(n\) dòng. Dòng thứ \(i\)
\((1\le i\le n)\) chứa \(n\) số nguyên \(c_{i1},c_{i2},\ldots,c_{in}\), trong
đó \(c_{ij}\) (\(1\le i,j\le n,\ 1\le c_{ij}\le n^2\)) là màu của bông hồng ở
hàng \(i\), cột \(j\).

Bảo đảm rằng có nhiều nhất \(10\) bộ dữ liệu có \(n>20\) và nhiều nhất \(4\)
bộ dữ liệu có \(n>100\).

\subsection*{Dữ liệu ra}

Với mỗi bộ dữ liệu, dòng đầu tiên in hai số nguyên \(d,m\)
\[
  0\le d\le n^2,\qquad 0\le m\le n^2,
\]
lần lượt là số loại màu khác nhau được dùng và số bông hồng cần sơn.

Trong \(m\) dòng tiếp theo, mỗi dòng in ba số nguyên \(i,j,c\)
\((1\le i,j\le n,\ 1\le c\le d)\), nghĩa là bạn sẽ sơn thân của bông hồng ở
hàng \(i\), cột \(j\) thành màu \(c\). Nếu thân của một bông hồng không được
sơn, nó được xem là có màu \(0\).

Miễn là số lượng loại màu được tối thiểu hóa, mọi cách sơn thân hoa hồng hợp
lệ đều được chấp nhận.

\subsection*{Ví dụ}

\paragraph{standard input}
\begin{verbatim}
3
3
1 1 1
1 1 1
1 1 1
3
1 2 1
2 1 2
1 2 1
3
1 2 3
4 5 6
7 8 9
\end{verbatim}

\paragraph{standard output}
\begin{verbatim}
2 6
1 2 2
1 3 1
2 1 1
2 3 2
3 1 2
3 2 1
1 4
1 3 1
3 1 1
2 3 1
3 2 1
0 0
\end{verbatim}

\section{Problem J. Unfair contest}

\paragraph{Tệp vào:} standard input
\paragraph{Tệp ra:} standard output
\paragraph{Giới hạn thời gian:} 1 second
\paragraph{Giới hạn bộ nhớ:} 256 megabytes

``Tôi sẽ chấm điểm công bằng. Chỉ là có thí sinh xứng đáng nhận một điểm số
công bằng hơn...''

gispzjz và zyb đang tham gia một cuộc thi, trong đó có \(n\) trọng tài cho
điểm họ, thường là theo màn trình diễn. Với mỗi thí sinh, mỗi trọng tài phải
nêu một số nguyên trong đoạn \([1,h]\) làm điểm, và điểm cuối cùng của thí
sinh là tổng tất cả điểm nhận được sau khi loại bỏ \(s\) điểm cao nhất và
\(t\) điểm thấp nhất.

Là một trong các trọng tài, bạn đã đặt cược vào gispzjz, nên bạn muốn anh ấy
thắng cuộc thi này, nhưng cũng không muốn việc đó trông quá lộ liễu. Giả sử
bạn biết \(n-1\) trọng tài còn lại đã cho gispzjz các điểm
\(a_1,\ldots,a_{n-1}\) và cho zyb các điểm \(b_1,\ldots,b_{n-1}\). Bạn cần
cho điểm \(a_n\) và \(b_n\) sao cho điểm cuối cùng của gispzjz lớn hơn hẳn
zyb. Nếu có thể làm được, bạn cũng cần tối thiểu hóa \(a_n-b_n\), với điều
kiện điểm cuối cùng của gispzjz lớn hơn hẳn zyb.

\subsection*{Dữ liệu vào}

Dòng đầu tiên chứa một số \(T\) (\(1\le T\le 12000\)), là số bộ dữ liệu.

Dòng đầu tiên của mỗi bộ dữ liệu chứa bốn số nguyên \(n,s,t,h\)
\[
  1\le n\le 10^5,\qquad 0\le s,t\le n-1,\qquad 1\le h\le 10^9,
\]
lần lượt là số trọng tài, số điểm cao nhất và thấp nhất cần loại bỏ, và phạm
vi điểm của trọng tài. Bảo đảm rằng \(s+t\le n-1\).

Sau đó là một dòng chứa \(n-1\) số nguyên \(a_1,\ldots,a_{n-1}\)
\((1\le a_i\le h)\), là các điểm đã được cho gispzjz.

Tiếp theo là một dòng chứa \(n-1\) số nguyên \(b_1,\ldots,b_{n-1}\)
\((1\le b_i\le h)\), là các điểm đã được cho zyb.

Bảo đảm rằng \(\sum n\le 10^6\) trên tất cả các bộ dữ liệu.

\subsection*{Dữ liệu ra}

Với mỗi bộ dữ liệu, nếu có thể làm cho điểm của gispzjz lớn hơn hẳn điểm của
zyb, hãy in giá trị nhỏ nhất của \(a_n-b_n\) trên một dòng. Ngược lại, in
\texttt{IMPOSSIBLE} trên một dòng.

\subsection*{Ví dụ}

\paragraph{standard input}
\begin{verbatim}
3
3 1 1 4
1 3
2 4
4 1 1 9
4 4 5
4 5 5
4 1 1 9
4 5 5
4 4 5
\end{verbatim}

\paragraph{standard output}
\begin{verbatim}
1
IMPOSSIBLE
-4
\end{verbatim}

\section{Problem K. ZYB's kingdom}

\paragraph{Tệp vào:} standard input
\paragraph{Tệp ra:} standard output
\paragraph{Giới hạn thời gian:} 4 seconds
\paragraph{Giới hạn bộ nhớ:} 512 megabytes

``Cha, thành phố của chúng ta và thành phố láng giềng đã nhiều năm không có
thu nhập. Chuyện gì đang xảy ra vậy?''

``Con à, ở thành phố láng giềng đang có dịch bệnh, con biết đấy...''

``Được rồi. Thế còn thành phố của chúng ta thì sao?''

``Ta thật muốn gõ vào đầu nhà vua khi ông ấy quyết định xây đất nước này như
một cái cây và biến thành phố của chúng ta thành một lá!''

Có \(n\) thành phố trong vương quốc của ZYB, được nối với nhau bằng \(n-1\)
con đường hai chiều. Ban đầu, thành phố thứ \(i\) có chỉ số thịnh vượng bằng
\(c_i\). Khi một giao dịch xảy ra giữa hai thành phố \(u\) và \(v\), thu nhập
của thành phố \(u\) tăng thêm \(c_v\), và thu nhập của thành phố \(v\) tăng
thêm \(c_u\).

Thỉnh thoảng, ZYB sẽ tổ chức một lễ hội giao thương để thúc đẩy kinh tế trong
vương quốc, và dĩ nhiên là để thu thêm thuế. Tuy nhiên, một số thành phố phải
bị đóng cửa và không được phép đi qua trong lễ hội do dịch bệnh. Nói hình
thức hơn, khi một lễ hội giao thương được tổ chức, một danh sách \(k\) thành
phố bị đóng cửa \(a_1,a_2,\ldots,a_k\) được thông báo. Với mọi cặp thành phố
\((u,v)\) (\(u<v\)), nếu đường đi duy nhất giữa chúng không chứa bất kỳ thành
phố nào trong \(k\) thành phố đó, thì một giao dịch xảy ra giữa chúng, tức là
thu nhập của thành phố \(u\) tăng thêm \(c_v\), và thu nhập của thành phố
\(v\) tăng thêm \(c_u\).

Ban đầu, thu nhập của mọi thành phố trong vương quốc của ZYB đều bằng \(0\).
Sau đó có \(q\) sự kiện xảy ra, mỗi sự kiện thuộc một trong các dạng sau:
\begin{enumerate}
  \item Tổ chức lễ hội giao thương. Cho danh sách \(k\) thành phố
  \(a_1,a_2,\ldots,a_k\). Với mọi cặp thành phố \((u,v)\) (\(u<v\)), nếu
  đường đi duy nhất giữa chúng không chứa bất kỳ thành phố nào trong danh sách
  này, thì một giao dịch xảy ra giữa chúng.
  \item Hỏi tổng thu nhập hiện tại của một thành phố \(v\).
\end{enumerate}

Bảo đảm rằng \(\sum n\le 10^6\), \(\sum q\le 10^6\) và \(\sum k\le 10^6\)
trên tất cả các bộ dữ liệu.

\subsection*{Dữ liệu vào}

Dòng đầu tiên chứa một số \(T\) (\(1\le T\le 20\)), là số bộ dữ liệu.

Dòng đầu tiên của mỗi bộ dữ liệu chứa hai số nguyên \(n,q\)
\((1\le n,q\le 2\cdot 10^5)\), lần lượt là số thành phố trong vương quốc của
ZYB và số sự kiện.

Sau đó là \(n-1\) dòng. Mỗi dòng chứa hai số nguyên \(u,v\), biểu thị một cạnh
giữa thành phố \(u\) và thành phố \(v\).

Tiếp theo là một dòng chứa \(n\) số nguyên \(c_1,c_2,\ldots,c_n\)
\((1\le c_i\le 10^6)\), là chỉ số thịnh vượng của từng thành phố.

Sau đó là \(q\) dòng, mỗi dòng mô tả một sự kiện thuộc một trong các dạng sau:
\begin{enumerate}
  \item \texttt{1 k} \((1\le k\le n)\), \(a_1,a_2,\ldots,a_k\)
  \((\forall\,1\le i\le k,\ 1\le a_i\le n\) và mọi \(a_i\) đôi một phân biệt),
  biểu thị một sự kiện loại một.
  \item \texttt{2 v} \((1\le v\le n)\), biểu thị một sự kiện loại hai hỏi thu
  nhập hiện tại của thành phố \(v\).
\end{enumerate}

\subsection*{Dữ liệu ra}

Với mỗi sự kiện loại hai, in một số trên một dòng, là đáp án.

\subsection*{Ví dụ}

\paragraph{standard input}
\begin{verbatim}
1
4 5
1 2
1 3
3 4
1 1 1 1
1 1 1
2 3
1 1 2
2 1
2 4
\end{verbatim}

\paragraph{standard output}
\begin{verbatim}
1
2
3
\end{verbatim}

\end{document}
